\scanpage{10}
\mspage{7}

\uline{Der}$_{2}$ Si $f:I\longrightarrow J$ est dans $\mathrm{Dia}$, si c'est une immersion \uline{fermée}, le biadjoint à droite $f^{!}_{\mathbb D}$ ou $f^{!}$ de $f^{*}$ existe. Si c'est une imm. ouverte, le \unsure{biadjoint} à gauche $f^{?}$ existe. Dans les deux cas, on veut que $f^{*}$ soit un foncteur-localisation, \deleted{\unsure{donc}} ce qui équivaut à dire que $f_{*}$ est pl. fid, ou que $f_{!}$ est pl. fidèle.

% slanted
\margin{Il n'y a lieu d'exiger les biadjoints que pour un dérivateur \uline{ponctué}. HOT n'y satisfait pas.}

% slanted
\margin{Il faut de plus \unsure{caractériser} l'image ess. de $f_{*}$ si $f$ immersion fermée (les $F$ tels que $g^{*}F\overset{\sim}{\longrightarrow}e_{\mathbb D(U)}$), et celle de $f_{!}$ si $f$ une imm. ouverte (les $F$ tels que $g^{*}(F)\simeq\unsure{\varphi|\emptyset}_{\mathbb D(Y)}$) où $g$ est l'immersion \unsure{compl\textsuperscript{re}|compl.}.}

\uline{NB} Dans le cas d'un dérivateur à gauche, on suppose \unsure{seult|seulemt} l'existence de $f^{!}$ si $f$ est une imm. fermée. Donc pour résumer, pour un tel dérivateur, on suppose l'existence de

\[
\left\{\begin{aligned}
&f_{*}\quad\text{quel que soit }f:I\unsure{-|\to}J\ \text{dans }\mathrm{Dia}\\
&f_{!}\quad\text{si }f\text{ est une imm. ouverte}\\
&f^{!}\quad\text{si }f\text{ est une imm. fermée.}
\end{aligned}\right.
\]

% slanted
\margin{$\swarrow$ cas ponctué seulement}

% slanted
\margin{\uline{$\mathrm{D\acute er}'_{2}$} : $\forall I,J\in\mathrm{Dia}$, on a $\mathbb D(I\amalg J)\underset{(\alpha^{*},\beta^{*})}{\overset{\approx}{\longrightarrow}}\mathbb D(I)\times\mathbb D(J)$, où $\alpha,\beta$ sont les \unsure{ap.|apl.} inj. can. de $I$ et $J$ dans $I\amalg J$.}

Dualement, pour un dérivateur à droite, on suppose l'existence de

\[
\left\{\begin{aligned}
&f_{!}\quad\text{quel que soit }f:I-J\ \text{dans }\mathrm{Dia}\\
&f_{*}\quad\text{si }f\text{ est une imm. fermée}\\
&f^{?}\quad\text{si }f\text{ est une imm. ouverte}
\end{aligned}\right.
\]

% slanted
\margin{$\swarrow$ cas ponctué}

Je voudrais formuler des axiomes astucieux, pour le cas d'un couple

\diagram{\begin{tikzcd}[ampersand replacement=\&]
\deleted{\unsure{X}}Y\arrow[r,hook,"i"] \& I \& U\arrow[l,hook',"j"']
\end{tikzcd}}

d'une imm. ouverte $j$ et d'une imm. fermée $i$ dans $\mathrm{Dia}$, compl. l'une de l'autre. D'abord une formulation si on suppose \uline{$\mathbb D$ ``ponctué''}, i.e. les catégories $\mathbb D(I)$ ponctuées, et les foncteurs $f^{*}$ \unsure{compatible|compatibles} avec les ponctuations. On veut ``interpréter'' le carré (1) page 2 :

% archive sheet number: 9
